\begin{lemma}\label{lem:06.2.4}
设 $O$ 是 $X$ 中的 $P_\Theta$-轨道，$C(w)$ 是 $O$ 中开的 Bruhat 胞腔。则对一切满足 $C(v)\not\subset O$ 的 $v\in W$ 与一切 $p\in\Z$，有 $R^pi_v^!(I(O,\lambda))=0$。若 $C(v)\subset O$，则当 $p\neq\ell(w)-\ell(v)$ 时 $R^pi_v^!(I(O,\lambda))=0$，且
$$R^{\ell(w)-\ell(v)}i_v^!(I(O,\lambda))=\calO_{C(v)}.$$
\end{lemma}

\begin{proof}
沿用前一个证明的记号。因 $O$ 仿射嵌入于 $X$，且 $I_O=R^0i_{O+}(\tau)$，第一个断言由基变换得出。记 $j_v$ 为 $C(v)$ 到 $O$ 的浸入。则由基变换可知，对任意 $p\in\Z$ 有
$$R^pi_v^!(I(O,\lambda))=R^pi_v^!(R^0i_{O+}(\tau))=R^pj_v^!(\tau).$$
另一方面，
$$R^pj_v^!(\tau)=L^{p-\dim O+\ell(v)}j_v^+(\tau),$$
它仅当 $p-\dim O+\ell(v)=0$，即 $p=\dim O-\ell(v)=\ell(w)-\ell(v)$ 时才非零；这是因为 $\tau$ 是联络。此外，$\tau$ 局部同构于 $\calO_O$，故 $L^0j_v^+(\tau)$ 是 $C(v)$ 上局部同构于 $\calO_{C(v)}$ 的 $N_0$-齐性联络，即 $L^0j_v^+(\tau)=\calO_{C(v)}$。
\end{proof}

由 1.8，对 $v\in W_\Theta$ 有 $\ell(vw)=\ell(w)-\ell(v)$，故
\begin{align*}
\nu(I(O,\lambda))&=\sum_{v\in W}\sum_{m\in\Z}\dim\bigl(R^mi_v^!(I(O,\lambda))\bigr)q^m\delta_v\\
&=\sum_{v\in W_\Theta}q^{\ell(v)}\delta_{vw}.
\end{align*}
如同 V.2.10 的证明，由此得到下述结果。

\begin{proposition}\label{prop:06.2.5}
设 $O$ 是 $X$ 中的 $P_\Theta$-轨道，$C(w)$ 是 $O$ 中开的 Bruhat 胞腔。则
$$\chi(I(O,\lambda))=\sum_{v\in W_\Theta}(-1)^{\ell(v)}\chi(I(vw,\lambda)).$$
\end{proposition}

特别地，若 $\Theta=\Pi$，则 $P_\Theta=G$ 传递地作用在 $X$ 上，且大胞腔就是附着于 $G$-轨道 $X$ 的 Bruhat 胞腔。于是有下述推论。

\begin{corollary}\label{cor:06.2.6}
$$\chi(\calO(\lambda+\rho))=\sum_{w\in W}(-1)^{\ell(w)}\chi(I(ww_0,\lambda)).$$
\end{corollary}

现在要描述如下最高权模：在 $\lambda$ 正则反支配时，它们经由范畴等价与标准 Harish-Chandra 层 $I(O,\lambda)$ 相对应。第一步是在最高权模范畴中构造一些对象。置
$$P^{++}(\Sigma_\Theta)=\{\nu\in P(\Sigma_\Theta)\mid \check\alpha(\nu)\in\Z_+,\ \alpha\in\Theta\}.$$
\begin{annote}
底本此处把条件中的变元印作 $\lambda$。但 $P^{++}(\Sigma_\Theta)$ 的定义中并无 $\lambda$，且下文 2.7 与 2.9 都要求 $\mu-\rho\in P^{++}(\Sigma_\Theta)$，即 $\check\alpha(\mu-\rho)\in\Z_+$。可见该条件是关于所论权 $\nu$ 的积分性。已改为 $\check\alpha(\nu)\in\Z_+$。[来源：底本 p.119，$P^{++}(\Sigma_\Theta)$ 的定义与 2.7 的用法]
\end{annote}

对任意 $\nu\in P^{++}(\Sigma_\Theta)$，用 $\bb_0$ 决定的特殊化，存在唯一的不可约有限维 $\lla_\Theta$-模 $V^\nu$，其最高权为 $\nu$。该模上有唯一的代数 $S_\Theta$-模作用，其微分给出 $s_\Theta$ 的作用。于是，若假定作用在 $\n_\Theta$ 与 $N_\Theta$ 上分别平凡，从而延拓到 $\p_\Theta$ 与 $Q_\Theta$，就可把 $V^\nu$ 看成 $(\p_\Theta,S_\Theta)$-模。对满足 $\mu-\rho\in P^{++}(\Sigma_\Theta)$ 的 $\mu\in\h^*$，定义广义 Verma 模
$$M_\Theta(\mu)=U(\g)\otimes_{U(\p_\Theta)}V^{\mu-\rho},$$
其中 $\g$ 的作用由第一个因子上的左乘给出，$Q_\Theta$ 的作用取为第一个因子上的伴随作用与第二个因子上的自然作用之张量积。显然，$M_\Theta(\mu)\in\mathcal{M}_{fg}(U(\g),Q_\Theta)$。

\begin{lemma}\label{lem:06.2.7}
设 $\mu\in\h^*$ 满足 $\mu-\rho\in P^{++}(\Sigma_\Theta)$。则 $M_\Theta(\mu)$ 是 Verma 模 $M(\mu)$ 的 $\p_\Theta$-有限商模中最大者。特别地，$M_\Theta(\mu)$ 是无穷小特征标为 $\chi_\mu$ 的最高权模。
\end{lemma}

\begin{proof}
$\p_\Theta$-模 $V^{\mu-\rho}$ 是 $\p_\Theta$-模 $U(\p_\Theta)\otimes_{U(\bb_0)}\C_{\mu-\rho}$ 的商。故 $M_\Theta(\mu)$ 是
\begin{align*}
U(\g)\otimes_{U(\p_\Theta)}\bigl(U(\p_\Theta)\otimes_{U(\bb_0)}\C_{\mu-\rho}\bigr)&=U(\g)\otimes_{U(\bb_0)}\C_{\mu-\rho}\\
&=M(\mu)
\end{align*}
的商。

设 $N$ 是 $M(\mu)$ 的 $\p_\Theta$-有限商。则 $N$ 含向量 $v$，它是生成元 $1\otimes1\in M(\mu)$ 的像。$v$ 是权为 $\mu-\rho$ 的权向量，且张成 $N$ 中的一维权子空间。设 $N'$ 是由 $v$ 生成的有限维 $\lla_\Theta$-子模。则 $N'$ 是若干不可约 $\lla_\Theta$-子模的直和，其中只有一个能含权 $\mu-\rho$ 对应的权子空间。这说明 $N'$ 其实不可约，且同构于 $V^{\mu-\rho}$。因此，$M(\mu)$ 到 $N$ 的投影通过 $M_\Theta(\mu)$ 分解。
\end{proof}

这说明，对 $\theta=W\cdot\mu$ 有 $M_\Theta(\mu)\in\mathcal{M}_{fg}(U_\theta,Q_\Theta)$。

已知 $\chi(M(\lambda))$（$\lambda\in\theta$）是 Grothendieck 群 $K(\mathcal{M}(U_\theta,N_0))$ 的一组基。因此，应当能把 $\chi(M_\Theta(\mu))$ 用 $\chi(M(\lambda))$（$\lambda\in\theta$）表示出来。

由 Poincaré--Birkhoff--Witt 定理，包络代数 $U(\g)$ 关于右乘是自由右 $U(\p_\Theta)$-模。这说明诱导函子 $V\mapsto U(\g)\otimes_{U(\p_\Theta)}V$ 从 $U(\p_\Theta)$-模范畴到 $U(\g)$-模范畴是正合的。若考虑 $s_\Theta$ 关于其 Borel 子代数 $\bb_0\cap\s_\Theta$ 的最高权模构成的范畴，则可通过把作用按某个线性形式延拓到 $\lla_\Theta$ 的中心、并平凡地延拓到 $\n_\Theta$，定义到 $U(\p_\Theta)$-模范畴的自然函子。该函子与诱导函子的复合给出从 $s_\Theta$ 的最高权模范畴到 $\g$ 的最高权模范畴的正合函子。因此，它给出相应 Grothendieck 群之间的态射。

下面需要一个简单的观察。

\begin{lemma}\label{lem:06.2.8}
设 $\nu\in P(\Sigma)$ 是支配权，$F^\nu$ 是最高权为 $\nu$ 的不可约有限维 $\g$-模。则在 $\g$ 的最高权模的 Grothendieck 群中有
$$\ch(F^\nu)=\sum_{w\in W}(-1)^{\ell(w)}\ch(M(w(\nu+\rho))).$$
\end{lemma}

\begin{proof}
$F^\nu$ 的最低权是 $w_0\nu$。由 2.6 有
$$\chi(\calO(w_0\nu))=\sum_{w\in W}(-1)^{\ell(w)}\chi(I(ww_0,w_0\nu-\rho)).$$
由范畴等价与 V.1.14，这说明
\begin{align*}
\ch(F^\nu)&=\sum_{w\in W}(-1)^{\ell(w)}\ch(I(w\nu-ww_0\rho))\\
&=\sum_{w\in W}(-1)^{\ell(w)}\ch(M(w(\nu+\rho))).
\end{align*}
\end{proof}

\begin{proposition}\label{prop:06.2.9}
设 $\mu\in\h^*$ 满足 $\mu-\rho\in P^{++}(\Sigma_\Theta)$。则
$$\ch(M_\Theta(\mu))=\sum_{v\in W_\Theta}(-1)^{\ell(v)}\ch(M(v\mu)).$$
\end{proposition}

\begin{proof}
$(W_\Theta,S_\Theta)$ 的长度函数是 $(W,S)$ 的长度函数 $\ell$ 在 $W_\Theta$ 上的限制。因此，把 2.6 用于 $\lla_\Theta$ 的表示 $V^{\mu-\rho}$，再结合前一段关于诱导函子所定义的 Grothendieck 群态射的观察，即得结论。
\end{proof}

置 $I_\Theta(\mu)=\widetilde{M_\Theta(\mu)}$。

\begin{theorem}\label{thm:06.2.10}
设 $\lambda\in\h^*$ 反支配。设 $O$ 是 $X$ 中的 $P_\Theta$-轨道，$w\in W$ 使 $C(w)$ 在 $O$ 中开。假定 $w\lambda\in P(\Sigma_\Theta)$，于是 $I(O,\lambda)$ 存在。则：
\begin{enumerate}[label=(\roman*)]
\item 对 $\alpha\in\Theta$ 有 $\check\alpha(w\lambda)\in\Z_+$；
\item 若对某个 $\alpha\in\Theta$ 有 $\check\alpha(w\lambda)=0$，则 $\Gamma(X,I(O,\lambda))=0$；
\item 若对一切 $\alpha\in\Theta$ 有 $\check\alpha(w\lambda)\neq0$，则 $w\lambda-\rho\in P^{++}(\Sigma_\Theta)$，且有
$$\Gamma(X,I(O,\lambda))\cong I_\Theta(w\lambda).$$
\end{enumerate}
一切满足 $\mu\in W\cdot\lambda$、$\mu-\rho\in P^{++}(\Sigma_\Theta)$ 的模 $I_\Theta(\mu)$，都由此得到。
\end{theorem}

\begin{proof}
因 $w\in{}^{\Theta}W$，(i) 成立。因 $\lambda$ 反支配时函子 $\Gamma$ 正合，由 2.5 与 V.1.14 有
\begin{align*}
\ch(\Gamma(X,I(O,\lambda)))&=\sum_{v\in W_\Theta}(-1)^{\ell(v)}\chi(\Gamma(X,I(vw,\lambda)))\\
&=\sum_{v\in W_\Theta}(-1)^{\ell(v)}\ch(I(vw\lambda)).
\end{align*}
因此，若对某个 $\alpha\in\Theta$ 有 $\check\alpha(w\lambda)=0$，则 $s_\alpha(w\lambda)=w\lambda$，从而 $\ch(\Gamma(X,I(O,\lambda)))=0$。这显然推出 (ii)。

若 (ii) 的假设不成立，则 $\check\alpha(w\lambda-\rho)=\check\alpha(w\lambda)-1\in\Z_+$，故 $w\lambda-\rho\in P^{++}(\Sigma_\Theta)$。不可约 $Q_\Theta$-齐性 $(\D_\lambda)_{i_O}$-联络 $\tau$ 在 $C(w)$ 上的限制等于 $\calO_{C(w)}$，故可把 $I(O,\lambda)$ 与 $I(w,\lambda)$ 的一个子模等同。因 $\lambda$ 反支配时 $\Gamma$ 正合，这说明 $\Gamma(X,I(O,\lambda))$ 是 $\Gamma(X,I(w,\lambda))$ 的 $U_\theta$-子模。由 V.1.14，$\Gamma(X,I(O,\lambda))$ 是 $I(w\lambda)$ 的 $\p_\Theta$-有限子模。把 2.7 的结论取对偶，可见 $\Gamma(X,I(O,\lambda))$ 是 $I_\Theta(w\lambda)$ 的子模。最后由上述计算与 2.9，得到
$$\ch(\Gamma(X,I(O,\lambda)))=\ch(I_\Theta(w\lambda)),$$
从而 $\Gamma(X,I(O,\lambda))\cong I_\Theta(w\lambda)$。

设 $\mu\in W\cdot\lambda$，$\mu-\rho\in P^{++}(\Sigma_\Theta)$。这说明对 $\alpha\in\Theta$ 有 $\check\alpha(\mu)\in\N$。取 $w\in W$，使对一切 $\beta\in w(\Sigma^+)$ 有 $-\check\beta(\mu)\in\Z_+$。则 $\lambda=w^{-1}\mu$ 反支配，且 $\Theta\subset-w(\Sigma^+)$，即 $w\in{}^{\Theta}W$。
\end{proof}

现讨论标准 Harish-Chandra 层 $I(O,\lambda)$ 的不可约性。为此先需要关于交织函子作用的一个结果。

\begin{lemma}\label{lem:06.2.11}
设 $O$ 是 $X$ 中的 $P_\Theta$-轨道，$C(w)$（$w\in{}^{\Theta}W$）是 $O$ 中开的 Bruhat 胞腔。设 $\alpha\in\Pi$ 满足 $ws_\alpha\in{}^{\Theta}W$、$\ell(ws_\alpha)=\ell(w)+1$，$O'$ 是对应的 $P_\Theta$-轨道。设 $w\lambda\in P(\Sigma_\Theta)$。则
$$I_{s_\alpha}(I(O,\lambda))=I(O',s_\alpha\lambda),$$
且 $L^{-1}I_{s_\alpha}(I(O,\lambda))=0$。
\end{lemma}

\begin{proof}
先注意 $(ws_\alpha)s_\alpha\lambda=w\lambda\in P(\Sigma_\Theta)$，故由 2.2，标准模 $I(O',s_\alpha\lambda)$ 存在。

设 $p_\alpha$ 是 $X$ 到型 $\alpha$ 的抛物子代数构成的广义旗簇 $X_\alpha$ 的投影。设 $Q$ 是 $X$ 中对应于 $C\in W_\Theta\backslash W$ 的 $P_\Theta$-轨道。则 $p_\alpha(Q)$ 是 $X_\alpha$ 中对应于 $W$ 中的双陪集 $W_\Theta\backslash W/W_{s_\alpha}$ 的 $P_\Theta$-轨道。因此，$p_\alpha^{-1}(p_\alpha(Q))$ 只有两种可能：
\begin{enumerate}[label=(\alph*)]
\item 若 $Cs_\alpha=C$，则是一个 $P_\Theta$-轨道；
\item 若 $Cs_\alpha\neq C$，则是两个 $P_\Theta$-轨道。
\end{enumerate}
按假设，$p_\alpha^{-1}(p_\alpha(O))=O\cup O'$，且 $\dim O'=\dim O+1$。于是
$$\dim p_\alpha(O)=\dim p_\alpha^{-1}(p_\alpha(O))-1=\dim O.$$

设 $x'\in O$。设 $P_{p_\alpha(x')}$ 是 $p_\alpha(x')$ 的稳定化子。则稳定化子 $P_\Theta\cap B_{x'}$ 与 $P_\Theta\cap P_{p_\alpha(x')}$ 维数相同；又它们都连通，故必相等。这说明 $p_\alpha^{-1}(p_\alpha(x'))\cap O$ 只含 $x'$。因此
$$p_\alpha^{-1}(p_\alpha(x'))=\{x'\}\cup\bigl(p_\alpha^{-1}(p_\alpha(x'))\cap O'\bigr),$$
且 $P_\Theta\cap P_{p_\alpha(x')}$ 传递地作用在 $p_\alpha^{-1}(p_\alpha(x'))\cap O'$ 上。于是最终得到：$P_\Theta\cap B_{x'}$ 传递地作用在 $p_\alpha^{-1}(p_\alpha(x'))\cap O'$ 上，即作用在一切相对于 $\bb_{x'}$ 处于相对位置 $s_\alpha$ 的 Borel 子代数构成的集上。

现沿用 L.3 的记号。设 $Z_{s_\alpha}\subset X\times X$ 是处于相对位置 $s_\alpha$ 的有序 Borel 子代数对构成的簇。因 $p_2^{-1}(O)$ 由一切满足 $x'\in O$、且 $\bb_x$ 相对于 $\bb_{x'}$ 处于相对位置 $s_\alpha$ 的有序对 $(x,x')\in X\times X$ 组成，由上述注记得知，$p_2^{-1}(O)$ 是维数为 $\dim O+1$ 的 $P_\Theta$-轨道。这说明 $p_1(p_2^{-1}(O))$ 也是 $P_\Theta$-轨道，故 $p_1(p_2^{-1}(O))=O'$。点 $x\in O'$ 在 $P_\Theta$ 中的稳定化子 $P_\Theta\cap B_x$ 包含 $(x,x')\in p_2^{-1}(O)$ 的稳定化子 $P_\Theta\cap B_x\cap B_{x'}$。由于轨道维数相同、稳定化子连通，可知 $p_1$ 诱导出 $p_2^{-1}(O)$ 到 $O'$ 的同构。

由 L.3.2 与 2.3，$p_2^{-1}(O)$ 仿射嵌入于 $Z_{s_\alpha}$。因此，记 $j$ 为 $p_2^{-1}(O)$ 到 $Z_{s_\alpha}$ 的仿射浸入，$q_2$ 为 $p_2$ 诱导的 $p_2^{-1}(O)$ 到 $O$ 的态射，则由基变换有
$$p_2^+(I(O,\lambda))=p_2^+(R^0i_{O+}(\tau))=R^0j_+(q_2^+(\tau)).$$
这说明
\begin{align*}
LI_{s_\alpha}(D(I(O,\lambda)))&=Rp_{1+}\bigl(T_{s_\alpha}\otimes_{\calO_{Z_{s_\alpha}}}D(R^0j_+(q_2^+(\tau)))\bigr)\\
&=R(p_1\circ j)_+(\tau'),
\end{align*}
其中 $\tau'$ 是 $p_2^{-1}(O)$ 上不可约的 $Q_\Theta$-等变联络。$p_1\circ j$ 的像等于 $O'$，且该映射是浸入。因此，$LI_{s_\alpha}(D(I(O,\lambda)))$ 等于 $D(V)$，其中 $V$ 是附着于 $O'$ 的标准 $\D_{s_\alpha\lambda}$-模。由 2.2，这个标准模就是 $I(O',s_\alpha\lambda)$。
\end{proof}

设 $w\in{}^{\Theta}W$。则 $\Sigma^+_\Theta\subset-w(\Sigma^+)$，从而 $-w^{-1}(\Sigma^+_\Theta)\subset\Sigma^+$。若 $w\lambda\in P(\Sigma_\Theta)$，则 $\Sigma_\Theta\subset\Sigma_{w\lambda}$，且 $w^{-1}(\Sigma_\Theta)\subset\Sigma_\lambda$。这说明
$$\Sigma^+_w\cap\Sigma_\lambda\supset-w^{-1}(\Sigma^+_\Theta).$$

\begin{theorem}\label{thm:06.2.12}
设 $O$ 是 $X$ 中的 $P_\Theta$-轨道，$C(w)$ 是 $O$ 中开的 Bruhat 胞腔。设 $w\lambda\in P(\Sigma_\Theta)$。则下列条件等价：
\begin{enumerate}[label=(\roman*)]
\item 标准模 $I(O,\lambda)$ 不可约；
\item $\Sigma^+_w\cap\Sigma_\lambda=-w^{-1}(\Sigma^+_\Theta)$。
\end{enumerate}
\end{theorem}

\begin{proof}
对 $\ell(w)$ 作归纳。若 $\ell(w)$ 最小，则 $w=w_\Theta$，$O$ 是 $X$ 中的闭 $P_\Theta$-轨道。此时 $I(O,\lambda)$ 总是不可约的。另一方面，$w_\Theta\lambda\in P(\Sigma_\Theta)$ 等价于 $\lambda\in P(\Sigma_\Theta)$，且
$$\Sigma^+_{w_\Theta}\cap\Sigma_\lambda=\Sigma^+_\Theta\cap\Sigma_\lambda=\Sigma^+_\Theta=-w_\Theta^{-1}(\Sigma^+_\Theta).$$

现设命题对 $\ell(w)<k$ 成立。设 $\ell(w)=k$。则由 1.7，可取 $\alpha\in\Pi$ 使 $w'=ws_\alpha\in{}^{\Theta}W$、$\ell(w')=k-1$。设 $O'$ 是对应于 $w'$ 的 $P_\Theta$-轨道。因 $w's_\alpha\lambda=w\lambda\in P(\Sigma_\Theta)$，标准模 $I(O',s_\alpha\lambda)$ 存在。于是由 2.11 有
$$I_{s_\alpha}(I(O',s_\alpha\lambda))=I(O,\lambda),$$
且 $L^{-1}I_{s_\alpha}(I(O',s_\alpha\lambda))=0$。

设 (ii) 成立。因 $\Sigma^+_w=s_\alpha(\Sigma^+_{w'})\cup\{\alpha\}$，可见
\begin{align*}
\Sigma^+_{w'}\cap\Sigma_{s_\alpha\lambda}&=s_\alpha\bigl((\Sigma^+_w-\{\alpha\})\cap\Sigma_\lambda\bigr)\\
&\subset-s_\alpha w^{-1}(\Sigma^+_\Theta)=-w'^{-1}(\Sigma^+_\Theta).
\end{align*}
由定理前的讨论，这说明上述包含是等式。于是由归纳假设，$I(O',s_\alpha\lambda)$ 不可约，且 $\alpha\notin\Sigma_\lambda$。这说明 $I_{s_\alpha}:\mathcal{M}_{qc}(\D_{s_\alpha\lambda})\to\mathcal{M}_{qc}(\D_\lambda)$ 是范畴等价，从而 $I(O,\lambda)$ 也不可约。

设 $I(O,\lambda)$ 不可约。若 $\check\alpha(\lambda)$ 不是整数，则 $I_{s_\alpha}:\mathcal{M}_{qc}(\D_{s_\alpha\lambda})\to\mathcal{M}_{qc}(\D_\lambda)$ 是范畴等价，故 $I(O',s_\alpha\lambda)$ 也不可约。由归纳假设，$\Sigma^+_{w'}\cap\Sigma_{s_\alpha\lambda}=-w'^{-1}(\Sigma^+_\Theta)$。因此
\begin{align*}
\Sigma^+_w\cap\Sigma_\lambda&=s_\alpha(\Sigma^+_{w'})\cap\Sigma_\lambda=s_\alpha(\Sigma^+_{w'}\cap\Sigma_{s_\alpha\lambda})\\
&=-s_\alpha w'^{-1}(\Sigma^+_\Theta)=-w^{-1}(\Sigma^+_\Theta).
\end{align*}

若 $\check\alpha(\lambda)$ 是整数，则由 L.7.3 与 2.11 有正合列
$$0\longrightarrow U^0\longrightarrow I(O,\lambda)\longrightarrow I(O',\lambda)\longrightarrow U^1\longrightarrow0,$$
其中中间的箭头非零。因 $I(O,\lambda)$ 的支集大于 $I(O',\lambda)$ 的支集，这与 $I(O,\lambda)$ 不可约的假设矛盾。
\end{proof}

下述结果由此推出。

\begin{theorem}\label{thm:06.2.13}
设 $\lambda\in P^{++}(\Sigma_\Theta)$ 正则。则下列条件等价：
\begin{enumerate}[label=(\roman*)]
\item $M_\Theta(\lambda)$ 不可约；
\item $\Sigma^+_\Theta=\{\alpha\in\Sigma^+\mid \check\alpha(\lambda)\in\N\}$。
\end{enumerate}
\end{theorem}

\begin{proof}
显然，可以改而考虑 $I_\Theta(\lambda)$。取 $w\in{}^{\Theta}W$ 使 $w^{-1}\lambda$ 反支配（由 2.10 的证明可知这样的 $w$ 存在）。设 $O$ 是附着于 $w$ 的左 $W_\Theta$-陪集的 $P_\Theta$-轨道。则 $\Gamma(X,I(O,w^{-1}\lambda))=I_\Theta(\lambda)$，且由范畴等价与 2.12，该模不可约当且仅当
$$\Sigma^+_w\cap\Sigma_{w^{-1}\lambda}=-w^{-1}(\Sigma^+_\Theta).$$
这等价于
$$\Sigma^+_{w^{-1}}\cap\Sigma_\lambda=\Sigma^+_\Theta.$$

设 $\beta\in\Sigma^+_{w^{-1}}\cap\Sigma_\lambda$。则 $\check\beta(\lambda)\in\Z$，且 $-\beta\in w(\Sigma^+)$，故 $-w^{-1}\beta\in\Sigma^+$，且
$$-\check\beta(\lambda)=-(w^{-1}\beta)^{\vee}(w^{-1}\lambda)\in-\N.$$
因此 $\check\beta(\lambda)\in\N$。反之，若对某个 $\beta\in\Sigma^+$ 有 $\check\beta(\lambda)\in\N$，则 $-\check\beta(\lambda)=-(w^{-1}\beta)^{\vee}(w^{-1}\lambda)\in-\N$；由于 $w^{-1}\lambda$ 反支配，必有 $-w^{-1}\beta\in\Sigma^+$。这说明 $\beta\in-w(\Sigma^+)$，从而 $\beta\in\Sigma^+_{w^{-1}}\cap\Sigma_\lambda$。因此
$$\Sigma^+_{w^{-1}}\cap\Sigma_\lambda=\{\beta\in\Sigma^+\mid \check\beta(\lambda)\in\N\}.$$
\end{proof}
